\documentclass[a4paper,10pt,twoside,openany]{article}

\usepackage[lang=hebrew]{maths}
\usepackage{hebrewdoc}
\usepackage{stylish}
\usepackage{lipsum}
\let\bs\blacksquare

\usepackage{siunitx,array}
\sisetup{table-format=-1.5, group-digits=false}
\newcolumntype{L}{>{\phantom{$-$}}l}       % prefix some whitespace
\newcommand\mcL[1]{\multicolumn{1}{L}{#1}} % handy shortcut macro

\setlength{\parindent}{0pt}

%%%%%%%%%%%%
% Styling %
%%%%%%%%%%%%

\usepackage{enumitem}

%%%%%%%%%%%%%
% Counters  %
%%%%%%%%%%%%%

\setcounter{section}{1}     
            
%BIBLIOGRAPHY
\usepackage[
backend=biber,
style=alphabetic,
]{biblatex}
\addbibresource{bibliography.bib} %Imports bibliography file

%%%%%%%%%%
% Title  %
%%%%%%%%%%
\title{
אלגברה ב' - טענה בנוגע לפולינום אופייני של כפל מטריצות
}
\date{}

\begin{document}
\maketitle

\begin{proposition}
תהיינה
$A,B \in \Mat_n\prs{\FF}$.
מתקיים
\begin{align*}
\text{.} p_{AB}\prs{x} = p_{BA}\prs{x}
\end{align*}
\end{proposition}

\begin{proof}
לפי הגדרת הפולינום האופייני, נרצה להראות כי
\begin{align*}
\text{.} \det\prs{x I_n - AB} = \det\prs{xI_n - BA}
\end{align*}
ניזכר כי
\begin{align}\label{eq:det_commutes}
det\prs{AB} = \det\prs{BA}
\end{align}
ונרצה להשתמש בכך, כיוון שמופיעה המטריצה
$AB$
באגף שמאל והמטריצה
$BA$
באגף ימין.

נוכל לנסות לכתוב את המטריצות שבשני האגפים ככפל של שתי מטריצות. למשל,
\begin{align*}
\text{.} x I_n - AB = x I_n \prs{I_n - \frac{1}{x} AB}
\end{align*}
אבל, לא נראה שהדטרמיננטה של המטריצה
$I_n - \frac{1}{x} AB$
קלה יותר לחישוב מזאת של
$x I_n - AB$.
לכן, ניעזר במשוואה
\eqref{eq:det_commutes}
בדרך אחרת, על ידי הסתכלות על מטריצות אחרות.

ניזכר כי הדטרמיננטה של מטריצה אלכסונית בלוקים היא מכפלת הדטרמיננטות של הבלוקים. נסמן
\begin{align*}
C &\coloneqq \pmat{AB & 0 \\ B & 0}
\end{align*}
שהדטרמיננטה שלה היא
$\det\prs{AB}$.
נרצה לכתוב דמיון בין
$C$
לבין מטריצה משולשת בלוקים שעל האלכסון שלה מופיע הבלוק
$BA$,
ואכן
\begin{align*}
\pmat{I_n & -A \\ 0 & I_n} C &= \pmat{0 & 0 \\ B & 0}
\\&= \pmat{0 & 0 \\ B & BA} \pmat{I_n & -A \\ 0 & I_n}
\end{align*}
ואז
$C$
דומה ל־%
\begin{align*}
D \coloneqq \pmat{0 & 0 \\ B & BA}
\end{align*}
שכן
$P^{-1} D P = C$
עבור
\[\text{.} P \coloneqq \pmat{I_n & - A \\ 0 & I_n}\]

כיוון שלמטריצות דומות יש אותו פולינום אופייני, נקבל כי
\begin{align*}
\text{.} \det\prs{x I_{2n} - C} = \det\prs{x I_{2n} - D}
\end{align*}
אגף שמאל שווה
\begin{align*}
\det\pmat{x I_n - AB & 0 \\ B & x I_n} = \det\prs{x I_n - AB} \det\prs{x I_n} = x^n \det\prs{x I_n - AB}
\end{align*}
ואילו אגף ימין שווה
\begin{align*}
\text{,} \det\pmat{x I_n & 0 \\ B & x I_n - BA} = \det\prs{x I_n} \det\prs{x I_n - BA} = x^n \det\prs{x I_n - BA}
\end{align*}
ולכן מתקיים
\begin{align*}
\text{,} \det\prs{x I_n - AB} = \det\prs{x I_n - BA}
\end{align*}
כנדרש.
\end{proof}

\end{document}